\documentclass[11pt]{article}
\usepackage[T1]{fontenc}
\usepackage[ngerman]{babel}
\usepackage{amsmath,amssymb,amsthm,mathtools}
\usepackage{geometry}
\geometry{margin=2.5cm}

\title{Quaternionische Rotationsdynamik und ein EABC-Modell \\ 
	\large Eine strukturierte algebraische Formulierung mit arithmetischer Interpretation}
\author{}
\date{}

\theoremstyle{definition}
\newtheorem{definition}{Definition}
\newtheorem{remark}{Bemerkung}
\theoremstyle{plain}
\newtheorem{lemma}{Lemma}
\newtheorem{proposition}{Proposition}

\begin{document}
	
	\maketitle
	
	\begin{abstract}
		Wir formulieren die klassische Rotationsdynamik vollständig in der Sprache der Quaternionen und isolieren die zugrunde liegende vierdimensionale Struktur. Diese wird anschließend als EABC-System interpretiert, wobei der skalare Anteil als Norm- bzw. Energiekanal und die imaginären Komponenten als gerichtete Dynamikachsen fungieren. Abschließend wird eine arithmetische Interpretation vorgeschlagen, die Quaternionennormen mit Vier-Quadrate-Darstellungen verknüpft. Die Arbeit trennt strikt zwischen bewiesenen Aussagen und heuristischen Erweiterungen.
	\end{abstract}
	
	\section{Quaternionen und Grundstruktur}
	
	\begin{definition}[Quaternion]
		Ein Quaternion ist ein Ausdruck
		\[
		q = q_0 + q_1 i + q_2 j + q_3 k,
		\]
		mit den Relationen
		\[
		i^2 = j^2 = k^2 = -1, \quad ij = k,\ jk = i,\ ki = j,
		\]
		sowie Antikommutativität
		\[
		ij = -ji.
		\]
	\end{definition}
	
	\begin{definition}[Norm]
		Die Norm eines Quaternionen ist
		\[
		\|q\|^2 = q_0^2 + q_1^2 + q_2^2 + q_3^2.
		\]
	\end{definition}
	
	\begin{lemma}[Multiplikationsstruktur]
		Für zwei reine Vektorquaternionen $a,b$ gilt
		\[
		ab = -\mathbf a \cdot \mathbf b + \mathbf a \times \mathbf b.
		\]
	\end{lemma}
	
	\begin{proof}
		Direkte Rechnung aus den Multiplikationsregeln.
	\end{proof}
	
	\section{Rotationen als Einheitsquaternionen}
	
	\begin{definition}[Rotationsquaternion]
		Sei $\mathbf n \in \mathbb{R}^3$ ein Einheitsvektor und $\theta \in \mathbb{R}$. Dann definiert
		\[
		q = \cos\frac{\theta}{2} + (n_x i + n_y j + n_z k)\sin\frac{\theta}{2}
		\]
		eine Rotation.
	\end{definition}
	
	\begin{proposition}[Konjugationswirkung]
		Für ein reines Vektorquaternion $v$ beschreibt
		\[
		v' = q v q^{-1}
		\]
		eine Drehung von $v$ um die Achse $\mathbf n$ mit Winkel $\theta$.
	\end{proposition}
	
	\section{Quaternionische Kinematik}
	
	\begin{definition}[Winkelgeschwindigkeit]
		Die Winkelgeschwindigkeit wird dargestellt durch
		\[
		\Omega = \omega_x i + \omega_y j + \omega_z k.
		\]
	\end{definition}
	
	\begin{proposition}[Bewegungsgleichung]
		Für eine zeitabhängige Rotation $q(t)$ gilt
		\[
		\dot q = \frac12 q \Omega.
		\]
	\end{proposition}
	
	\begin{remark}
		Die Nichtkommutativität der Quaternionen entspricht der Nichtvertauschbarkeit von Rotationen.
	\end{remark}
	
	\section{Dynamik}
	
	\begin{definition}[Drehimpuls]
		\[
		\Lambda = \mathcal I(\Omega),
		\]
		wobei $\mathcal I$ ein linearer Operator (Trägheitstensor) ist.
	\end{definition}
	
	\begin{definition}[Drehmoment]
		\[
		T = \frac{d\Lambda}{dt}.
		\]
	\end{definition}
	
	\begin{definition}[Energie]
		\[
		K = \frac12 \langle \Omega, \Lambda \rangle.
		\]
	\end{definition}
	
	\begin{definition}[Leistung]
		\[
		P = \langle T, \Omega \rangle.
		\]
	\end{definition}
	
	\section{EABC-Struktur}
	
	\begin{definition}[EABC-Zuordnung]
		Wir definieren eine Identifikation
		\[
		E \leftrightarrow 1,\quad A \leftrightarrow i,\quad B \leftrightarrow j,\quad C \leftrightarrow k.
		\]
	\end{definition}
	
	\begin{definition}[EABC-Zustand]
		Ein Zustand ist gegeben durch
		\[
		Q = E + Ai + Bj + Ck.
		\]
	\end{definition}
	
	\begin{proposition}[Normerhaltung]
		Unter der Transformation
		\[
		Q \mapsto q Q q^{-1}
		\]
		bleibt die Norm
		\[
		N(Q) = E^2 + A^2 + B^2 + C^2
		\]
		erhalten.
	\end{proposition}
	
	\begin{remark}
		Dies entspricht einer Rotation im dreidimensionalen Unterraum $(A,B,C)$ bei fixem Normwert.
	\end{remark}
	
	\section{Arithmetische Interpretation}
	
	\begin{definition}[Arithmetischer Zustand]
		Sei
		\[
		Q_N = E_N + A_N i + B_N j + C_N k,
		\]
		mit
		\[
		N = E_N^2 + A_N^2 + B_N^2 + C_N^2.
		\]
	\end{definition}
	
	\begin{proposition}
		Jede solche Darstellung entspricht einer Vier-Quadrate-Zerlegung von $N$.
	\end{proposition}
	
	\begin{remark}
		Dies ist eine direkte Verbindung zum Satz von Lagrange.
	\end{remark}
	
	\section{Heuristische Erweiterung: Dynamik im EABC-Raum}
	
	\begin{definition}[Quaternionische Dynamik]
		\[
		\dot Q = \frac12 Q \Omega.
		\]
	\end{definition}
	
	\begin{remark}[Interpretation]
		\begin{itemize}
			\item $E$: skalarer Norm- bzw. Energieanteil
			\item $A,B,C$: gerichtete Dynamikachsen
		\end{itemize}
	\end{remark}
	
	\begin{remark}[Heuristik]
		Die Größen
		\[
		\Omega,\ \Lambda,\ T
		\]
		können als ``Ströme'' zwischen den Kanälen $A,B,C$ interpretiert werden.
	\end{remark}
	
	\section{Symmetrisches Beispiel}
	
	\begin{proposition}
		Für $\mathbf n = (1,1,1)/\sqrt3$ und $\theta = 2\pi/3$ gilt
		\[
		q = \frac{1+i+j+k}{2}.
		\]
	\end{proposition}
	
	\begin{remark}
		Dieses Quaternion erzeugt eine zyklische Permutation der Achsen und besitzt tetraedrische Symmetrie.
	\end{remark}
	
	\section{Diskussion}
	
	Die quaternionische Darstellung zeigt:
	
	\begin{itemize}
		\item Energie ist eine skalare Projektion
		\item Dynamik lebt im imaginären Teil
		\item Nichtkommutativität erzeugt echte Kopplungseffekte
	\end{itemize}
	
	Die EABC-Interpretation ist konsistent mit:
	
	\begin{itemize}
		\item Vier-Kanal-Struktur
		\item Normerhaltung
		\item Rotationssymmetrie
	\end{itemize}
	
	\section{Ausblick}
	
	Die hier dargestellte Struktur legt nahe:
	
	\begin{itemize}
		\item Dynamische Systeme auf Hurwitz-Quaternionen
		\item Diskrete Versionen mit ganzzahligen Zuständen
		\item Anwendungen auf arithmetische Strukturen (Primzahlklassen)
	\end{itemize}
	
	Eine vollständige mathematische Theorie dieser Erweiterung bleibt Gegenstand zukünftiger Arbeit.
	
	\section*{Quaternionic Symmetries of Prime Quadruplets and Residue Dynamics}
	
	\subsection*{Abstract}
	We investigate a quaternionic embedding of prime quadruplets and study their transformation under rotations in $\mathbb{H}$. While exact quadruplet structure is not preserved under such transformations combined with projection onto primes, we show that residue class patterns modulo $12$ exhibit statistical stability. This suggests an underlying geometric symmetry in the distribution of prime quadruplets.
	
	\subsection*{1. Quaternionic Embedding}
	Define the embedding
	[
	Q(p) = p + (p+2)i + (p+6)j + (p+8)k
	]
	for prime quadruplets.
	
	\subsection*{2. Transformation}
	Let
	[
	u = \frac{1+i+j+k}{2}
	]
	and define
	[
	T_u(Q) = u(Q - Q_{\text{mid}})u^{-1} + Q_{\text{mid}}.
	]
	
	\subsection*{3. Projection}
	Define a projection
	[
	\Pi(Q)_k = \text{nearest prime to } Q_k.
	]
	
	\subsection*{4. Main Operator}
	[
	\mathcal{T} = \Pi \circ T_u.
	]
	
	\subsection*{5. Structural Result}
	The operator $\mathcal{T}$ does not preserve prime quadruplets, but preserves residue class structure modulo $12$ in a statistical sense.
	
	\subsection*{6. Interpretation}
	This induces a dynamical system on residue classes corresponding to a permutation of $(5,7,11)$, compatible with quaternionic rotations.
	
	\subsection*{7. Outlook}
	Possible connections to spectral operators and SU(2) symmetry remain open.
	
	
\end{document}